% Lösungen Einheit 1
\einheitenkopf*{Einheit 1 von 5 · Lineares und exponentielles Wachstum unterscheiden}

\erg{10}{a) gleicher Betrag \quad b) gleicher Prozentsatz \quad c) gleicher Betrag \quad d) gleicher Prozentsatz \quad e) gleicher Prozentsatz}
\erg{11}{a) gleicher Quotient \quad b) gleiche Differenz \quad c) gleicher Quotient \quad d) gleiche Differenz \quad e) gleicher Quotient}
\erg{12}{a) Graph I bei $0$ (Ursprung) \quad b) Graph II bei $1{,}2$ \quad c) Graph III bei $2$ \quad d) Graph I}
\erg{13}{a) $7;\ 7;\ 7$ – ja \quad b) $10;\ 10;\ 10$ – ja \quad c) $3;\ 4;\ 5$ – nein \quad d) $15;\ 15;\ 15$ – ja \quad e) $1{,}1;\ 1{,}1;\ 1{,}1$ – ja \quad f) exponentiell: die Quotienten sind immer $1{,}05$ (die Differenzen $40;\ 42;\ 44{,}1$ nicht) \quad g) exponentiell: derselbe Prozentsatz heißt jeden Monat mal $1{,}15$ \quad h) exponentiell, nimmt ab: die Quotienten sind immer $0{,}9$ \quad i) $40;\ 42;\ 44{,}1$ – die $5\,\%$ werden jedes Jahr vom größeren Wert genommen \quad j) Der Teich ist begrenzt: Ist er zugewachsen, kann die Fläche nicht weiter wachsen. \quad k) $252\,000 : 240\,000 = 1{,}05$; $264\,600 : 252\,000 = 1{,}05$; $277\,830 : 264\,600 = 1{,}05$ – gleiche Quotienten, also exponentiell (Differenzen $12\,000;\ 12\,600;\ 13\,230$ sind nicht gleich)}
\erg{14}{a) Gerade \quad b) Kurve \quad c) Gerade \quad d) Kurve \quad e) I und IV (beginnen im Ursprung) \quad f) II}
\erg{15}{a) Eine Gerade fällt jedes Jahr um denselben Betrag ($300\,€$); bei $15\,\%$ wird der Verlust jedes Jahr kleiner. \quad b) Graph III beginnt bei $2\,000\,€$, nicht bei $3\,200\,€$. \quad c) Graph II beginnt bei $3\,200\,€$, fällt und wird immer flacher, weil jedes Jahr $15\,\%$ vom kleineren Wert abgehen. \quad d) Graph III passt (beginnt bei $40$ Tausend, steigt immer steiler). Graph I ist eine Gerade: jeden Monat gleicher Zuwachs, nicht gleicher Prozentsatz. Graph II beginnt im Ursprung: am Anfang $0$ statt $40\,000$ Abonnenten.}
\erg{16}{a)–d) Punkte $(0\,|\,10)$, $(1\,|\,20)$, $(2\,|\,40)$, $(3\,|\,80)$ \quad e) $(2{,}5\,|\,57)$ \quad f) Kurve durch alle Punkte, die immer steiler wird}
\erg{17}{a) z.\,B. waagerecht 1 Kästchen $= 1$ Jahr, senkrecht 1 Kästchen $= 100\,€$; Punkte $(0\,|\,500)$, $(2\,|\,605)$, $(4\,|\,732)$, $(6\,|\,886)$, $(8\,|\,1\,072)$ \quad b) z.\,B. waagerecht 2 Kästchen $= 1$ m, senkrecht 1 Kästchen $= 100$ Lux; Punkte $(0\,|\,800)$, $(1\,|\,520)$, $(2\,|\,338)$, $(3\,|\,220)$; fallende Kurve, die flacher wird}
\erg{18}{a) Fehler: Gleicher Prozentsatz heißt nicht gleicher Betrag. Die Zuwächse sind $60\,€$, $67{,}20\,€$, $75{,}26\,€$, die Quotienten immer $1{,}12$ – exponentielles Wachstum. \quad b) exponentiell (Abnahme): Quotienten immer $0{,}98$; Differenzen $700;\ 686;\ 672$ nicht gleich}
\erg{19}{a) Nein: Die $2\,\%$ beziehen sich jedes Jahr auf das größere Guthaben, die Zinsen werden größer. \quad b) kleiner: $10\,\%$ vom jeweils kleineren Wert sind weniger Euro. \quad c) Nein: Zur Zeit $0$ hat der Bestand den Wert $500$, der Graph beginnt also bei $500$ auf der senkrechten Achse.}
\erg{20}{a) Plan A linear (jedes Jahr $+150$), Plan B exponentiell (jedes Jahr mal $1{,}1$) \quad b) ab Jahr 6 ($2\,126 > 2\,100$) \quad c) mit Plan B in Jahr 7 ($2\,338 > 2\,300$); mit Plan A erst in Jahr 8}
